\chapter{Stratégie technologique}\label{ch:fm:technology-strategy}

Les dépenses de communication et de traitement des données de Virtu Financial s'élevaient à 209~millions de dollars en 2019 et à 249~millions de dollars en 2025. Dans l'intervalle, son chiffre d'affaires est passé de
1.5~milliards de dollars à 3.2~milliards, est redescendu à 2.3~milliards puis est remonté à 3.6~milliards~; la dépense a augmenté d'un cinquième et n'a baissé qu'une seule fois, de moins de 1~\%. Une société de trading ne
choisit pas chaque année s'il faut être rapide. Le \omterm{def:fm:technology-strategy:tier}{palier de latence} qu'elle paie est une décision sur les activités auxquelles elle entend se consacrer, prise une fois pour toutes et payée chaque
année, quelle que soit la conjoncture.

\section{L'architecture comme stratégie}

Le livre~13 a construit la pile à faible latence et le livre~14 les réseaux qui la sous-tendent. Ce chapitre porte sur la décision qui se situe au-dessus des deux~: quelle vitesse acheter, pour
quelles stratégies, et ce que le même budget pourrait financer d'autre. La décision est stratégique parce que les coûts sont fixes (la \omterm{def:fm:the-economics-of-a-trading-firm:fixed}{base de coûts fixes} du chapitre~1~; la \cref{fig:fm:technology-strategy:virtu}) et
que les bénéfices dépendent de ce que les concurrents achètent.

\begin{omfigure}
\begin{tikzpicture}
\begin{axis}[width=10.5cm,height=5.2cm,ymode=log,xlabel={\small année},ylabel={\small millions de dollars (échelle logarithmique)},xmin=2018.6,xmax=2025.4,ymin=150,
  ymax=5000,xtick={2019,2020,2021,2022,2023,2024,2025},xticklabel style={/pgf/number format/1000 sep=},ytick={200,500,1000,2000,4000},
  yticklabels={200,500,1000,2000,4000},tick label style={font=\footnotesize},grid=major,grid style={gray!20},table/col sep=comma,
  legend cell align=left,legend style={font=\scriptsize,at={(0.5,-0.3)},anchor=north,/tikz/every even column/.append style={column sep=8pt}},legend columns=2]
\addplot[omDef,thick,mark=*,mark size=1.3pt] table[x=year,y=revenue]{figdata/desk/19-technology-strategy/virtu.csv};
\addplot[omProp,thick,mark=square*,mark size=1.3pt] table[x=year,y=comm_data]{figdata/desk/19-technology-strategy/virtu.csv};
\legend{chiffre d'affaires total,communication et traitement des données}
\end{axis}
\end{tikzpicture}

\omcaption{Chiffre d'affaires total de Virtu Financial et sa dépense de communication et de traitement des données, 2019--2025. Le chiffre d'affaires a plus que doublé puis s'est replié~;
la dépense a augmenté régulièrement. Source~: Forms 10-K (le registre du chapitre)~; données dans \texttt{data/desk/filings\_virtu.csv}.}\label{fig:fm:technology-strategy:virtu}
\end{omfigure}

\begin{definition}[Palier de latence]\label{def:fm:technology-strategy:tier}
Un \emph{palier de latence}\index{palier de latence} est un niveau d'infrastructure de trading défini par la latence qu'il atteint et par le coût fixe qu'il entraîne~:
par exemple, des connexions partagées et un trading logiciel~; des serveurs hébergés en colocation avec des liaisons directes à faible latence~; ou un trading matériel en colocation avec
des routes hertziennes dédiées entre les plateformes.
\end{definition}

Les trois paliers du chapitre sont chiffrés à partir des frais de plateforme que le livre~14 a relevés dans les grilles tarifaires publiées, auxquels s'ajoutent le matériel, les routes et les équipes de la société,
qui sont des données d'entrée (l'\cref{dat:fm:technology-strategy:fees})~: 0.29~million de dollars par an pour le palier~1, 1.35~million de dollars pour le palier~2 et 4.45~millions de dollars pour le palier~3,
par stratégie. Les frais de plateforme ne représentent qu'une faible part~: 36\,000, 450\,000 et 946\,000 dollars par an.

\begin{dated}{2026-09}{Frais de plateforme derrière les paliers}\label{dat:fm:technology-strategy:fees}
D'après les grilles tarifaires relevées dans le livre~14 (chapitre~9)~: dans le centre de données NY11-4 de Nasdaq, une baie à très haute densité facturée 7\,230~dollars par mois,
avec des frais d'installation de 5\,940~dollars pour la baie, 4\,560~dollars pour l'alimentation de phase~3 et 5\,260~dollars pour une PDU de phase~3 (SEC Release 34-101267, octobre
2024)~; chez MIAX Pearl Options, une connexion à très faible latence de 10~Gb facturée 15\,000~dollars par mois, une connexion de reprise après sinistre de 10~Gb à 3\,500~dollars et une connexion de 1~Gb
à 1\,500~dollars (Federal Register 2026-10452, mai 2026). Le palier~1 compte deux connexions de 1~Gb~; le palier~2 une baie et deux connexions de 10~Gb~; le palier~3
deux baies, quatre connexions de 10~Gb et une liaison de reprise après sinistre, les frais ponctuels étant répartis sur 36~mois. Les grilles tarifaires évoluent par dépôt de documents.
\end{dated}

\section{Ce que la vitesse rapporte, stratégie par stratégie}

La vitesse rapporte de l'argent dans les courses~: deux firmes ou plus réagissent au même événement public, et seule la première obtient la transaction. Aquilina, Budish et O'Neill
(2022) ont mesuré directement ces courses dans les données de messages de la bourse pour les actions du FTSE~100~: environ une course par minute et par symbole, d'une durée de 5 à 10
millionièmes de seconde dans le cas le plus fréquent, soit environ un cinquième du volume échangé, avec six firmes qui concentrent plus de 80~\% des gains et des pertes de ces courses. Les courses
représentaient à peu près un tiers du coût de la liquidité. Pour une stratégie, la question est de savoir quelle part de son chiffre d'affaires est décidée par les courses.

\begin{definition}[Course aux armements de la latence]\label{def:fm:technology-strategy:race}
Une \emph{course aux armements de la latence}\index{course aux armements de la latence} est une concurrence dans laquelle des firmes achètent de la vitesse pour gagner des courses les unes contre les autres, de sorte que chaque firme
voit ses investissements réduire le rendement des autres, et la dépense d'équilibre en vitesse dépasse ce que les firmes gagnent ensemble.
\end{definition}

Le modèle de course du chapitre répartit le fonds de revenus $V$ d'une stratégie, partagé entre $n$ concurrents, en une part de course $\alpha$ et le reste. Les courses reviennent aux
firmes du palier le plus rapide présent, partagées à parts égales~; le reste est partagé à parts égales entre les $n$. Une firme gagne
\[
V\Bigl(\alpha\,\frac{\mathbf 1\{\text{dans le palier supérieur}\}}{m}+\frac{1-\alpha}{n}\Bigr)-c_{\text{palier}},
\]
où $m$ est le nombre de firmes dans le palier supérieur (le \cref{lst:fm:technology-strategy:race}). Cinq stratégies diffèrent par $V$ et par $\alpha$~: le market making sur contrats à terme
(60~millions de dollars, 0.8), l'arbitrage sur ETF (30~millions de dollars, 0.9), l'arbitrage statistique sur actions à l'échelle de la minute (20~millions de dollars, 0.1), le market making sur options
(40~millions de dollars, 0.5) et le basis trade crypto (10~millions de dollars, 0.3). Les valeurs sont illustratives~; les parts de course sont fixées au regard du
constat selon lequel les courses représentent un cinquième du volume et un tiers du coût de la liquidité sur un marché où les teneurs de marché dominent les deux.

\omcode{code/firm/techtier/firm_techtier.py}{52}{76}{Le modèle de course~: le revenu d'une firme selon son palier et celui de ses rivaux, son meilleur palier et la course aux armements
résolue par réponses optimales.}\label{lst:fm:technology-strategy:race}

\section{Tutoriel~: le palier le plus rapide}\label{tut:fm:technology-strategy:tut}

\begin{tutorial}
\textbf{Objectif.} Choisir un \omterm{def:fm:technology-strategy:tier}{palier de latence} pour chacune des cinq stratégies face à des rivaux donnés, puis laisser chaque firme choisir et trouver
l'équilibre de la course aux armements. \textbf{État final~:} le graphique des profits par palier (la \cref{fig:fm:technology-strategy:profits}) et la dépense d'équilibre comparée à l'absence de mise à niveau
(la \cref{fig:fm:technology-strategy:race}).
\begin{tutsteps}
\item \textbf{Les paliers.} \texttt{fm\_tech.tiers()} construit les trois paliers à partir des grilles \texttt{firm.colobill} du livre~14 et des données d'entrée indiquées.
\item \textbf{Les rivaux.} Chaque stratégie compte quatre rivaux à des paliers donnés~: pour le market making sur contrats à terme, deux au palier~3, un au palier~2 et un au palier~1.
\item \textbf{Le choix.} \texttt{firm.techtier.best\_tier} compare les profits des trois paliers pour chaque stratégie.
\item \textbf{La course.} \texttt{firm.techtier.arms\_race} fait démarrer cinq firmes identiques au palier~1 et laisse chacune choisir à son tour son meilleur palier jusqu'à
ce qu'aucune firme ne change~; \texttt{waste} compare la dépense d'équilibre à celle où tout le monde reste au palier~1.
\end{tutsteps}
\end{tutorial}

\begin{omfigure}
\begin{tikzpicture}
\begin{axis}[width=10.5cm,height=5.4cm,ybar,bar width=7pt,ylabel={\small profit (millions de dollars par an)},ymin=-2,ymax=22,xtick={0,1,2,3,4},
  xticklabels={{MM futures},{arb ETF},{arb stat},{MM options},{crypto}},xmin=-0.6,xmax=4.6,
  tick label style={font=\footnotesize},grid=major,grid style={gray!20},table/col sep=comma,
  legend cell align=left,legend style={font=\scriptsize,at={(0.5,-0.25)},anchor=north,/tikz/every even column/.append style={column sep=8pt}},legend columns=3]
\addplot[fill=gray!40,draw=gray] table[x=pos,y=t1]{figdata/desk/19-technology-strategy/profits.csv};
\addplot[fill=omThm!50,draw=omThm] table[x=pos,y=t2]{figdata/desk/19-technology-strategy/profits.csv};
\addplot[fill=omDef!55,draw=omDef] table[x=pos,y=t3]{figdata/desk/19-technology-strategy/profits.csv};
\legend{palier 1,palier 2,palier 3}
\end{axis}
\end{tikzpicture}

\omcaption{Profit par \omterm{def:fm:technology-strategy:tier}{palier de latence} pour chacune des cinq stratégies, face à des rivaux donnés. La vitesse paie là où les courses décident de l'essentiel du revenu et où le fonds
est important~; le palier~3 rapporte moins pour le book d'arbitrage statistique et fait perdre de l'argent sur le basis trade crypto. Données~: \texttt{fm\_tech.choices}.}\label{fig:fm:technology-strategy:profits}
\end{omfigure}

Face à ses rivaux, la firme devrait acheter le palier~3 pour le market making sur contrats à terme (profit de 13.95~millions de dollars par an, contre 2.11~millions au palier~1), l'arbitrage sur ETF
(9.65~millions de dollars contre 0.31~million) et le market making sur options (19.55~millions de dollars contre 3.71~millions)~; le palier~2 pour le basis trade crypto
(3.05~millions de dollars), où aucun rival n'en a davantage~; et le palier~1 pour l'arbitrage statistique à l'échelle de la minute (3.31~millions de dollars), qui perd 2.16~millions de dollars par an
en achetant le palier~3. La même firme est une entreprise de vitesse dans trois stratégies et non dans deux, et son architecture doit le refléter.

\begin{proposition}[Combien de firmes courent]\label{prop:fm:technology-strategy:count}
Avec deux paliers dont les coûts diffèrent de $\Delta c$, $n$ firmes identiques et le modèle de course, une firme a intérêt à rejoindre $m-1$ autres dans le palier rapide si
et seulement si $\alpha V/m\ge\Delta c$. À l'équilibre, le nombre de firmes dans le palier rapide est $m^*=\min(n,\lfloor\alpha V/\Delta c\rfloor)$, leur
dépense totale en vitesse dépasse de $m^*\Delta c$ le niveau où toutes sont lentes, et le fonds partagé par l'industrie reste inchangé.
\end{proposition}

\begin{proof}
Passer du palier lent au palier rapide fait passer le revenu de course d'une firme de 0 à $\alpha V/m$ et laisse sa part du reste à $(1-\alpha)V/n$~;
le coût augmente de $\Delta c$. Les firmes rejoignent le palier tant que $\alpha V/m\ge\Delta c$, de sorte que la dernière à entrer est la $\lfloor\alpha V/\Delta c\rfloor$-ième, et une firme
du palier rapide ne le quitte pas, car le quitter fait perdre $\alpha V/m^*\ge\Delta c$. Les revenus font la somme $V$ quels que soient les paliers.
\end{proof}

Avec trois paliers, le raisonnement s'applique étape par étape~: pour le market making sur contrats à terme, $\alpha V=48$ et le passage du palier~1 au palier~3 coûte 4.16~millions de dollars,
donc jusqu'à onze firmes pourraient courir~; avec cinq, les cinq courent. Le market making sur options donne $20/4.16$, soit quatre firmes~; le basis trade crypto, $3/1.06$ pour
le passage au palier~2, soit deux firmes~; l'arbitrage statistique, $2/1.06$, soit une seule.

\begin{omfigure}
\begin{tikzpicture}
\begin{axis}[width=10.5cm,height=5.2cm,ybar,bar width=10pt,ylabel={\small dépense en vitesse (millions de \$ par an)},ymin=0,ymax=25,xtick={0,1,2,3,4},
  xticklabels={{MM futures},{arb ETF},{arb stat},{MM options},{crypto}},xmin=-0.6,xmax=4.6,
  tick label style={font=\footnotesize},grid=major,grid style={gray!20},table/col sep=comma,
  legend cell align=left,legend style={font=\scriptsize,at={(0.5,-0.25)},anchor=north,/tikz/every even column/.append style={column sep=8pt}},legend columns=2]
\addplot[fill=omDef!55,draw=omDef] table[x=pos,y=spend]{figdata/desk/19-technology-strategy/race.csv};
\addplot[fill=gray!40,draw=gray] table[x=pos,y=baseline]{figdata/desk/19-technology-strategy/race.csv};
\legend{équilibre de la course aux armements,personne ne se met à niveau}
\end{axis}
\end{tikzpicture}

\omcaption{Dépense du secteur en vitesse, pour cinq firmes identiques dans chaque stratégie, à l'équilibre de la course aux armements et si personne ne s'était mis à niveau. Le
fonds de revenus est le même dans les deux cas~; la différence est dépensée pour le redistribuer. Données~: \texttt{fm\_tech.races}.}\label{fig:fm:technology-strategy:race}
\end{omfigure}

\section{Paliers de latence et course aux armements}

L'équilibre dépense 22.2~millions de dollars par an en vitesse pour le market making sur contrats à terme et l'arbitrage sur ETF, contre 1.4~million de dollars si personne ne s'était mis à niveau~: 93.6~\%
de la dépense n'achète rien pour le secteur, puisque le fonds est le même. Le profit du secteur en market making sur contrats à terme tombe de 58.6~millions de dollars
à 37.8~millions de dollars~; en arbitrage sur ETF, de 28.6~millions de dollars à 7.8~millions, soit environ un quart de ce qu'il serait. Le market making sur options gaspille 92.1~\%, le
basis trade crypto 59.8~\% et l'arbitrage statistique 42.7~\%. Pour les cinq, le secteur dépense 68.6~millions de dollars là où 7.2~millions suffiraient~: 89.6~\% de
sa dépense en vitesse est du gaspillage au sens du modèle. C'est l'argument que Budish, Cramton et Shim (2015) avancent pour les marchés en général~: la concurrence ne réduit pas
les opportunités d'arbitrage, elle élève la vitesse nécessaire pour les saisir.

Pour une firme isolée, la conclusion est l'inverse de celle du secteur~: dans une course que d'autres mènent, ne pas acheter de vitesse revient à perdre les revenus de la course.
La décision de la firme porte sur les courses où elle se lance et celles où elle ne se lance pas, et, lorsque le fonds est petit ou les courses peu nombreuses, sur le fait de se retirer.

\section{Équipes plateforme et équipes produit}

\begin{definition}[Équipe plateforme, équipe produit]\label{def:fm:technology-strategy:teams}
Une \emph{équipe produit}\index{équipe produit} construit et exploite ce qu'une activité utilise pour gagner de l'argent, comme le système de trading d'une stratégie, et est
mesurée sur les résultats de cette activité. Une \emph{équipe plateforme}\index{équipe plateforme} construit et exploite des capacités partagées que les équipes produit utilisent
sous forme de service, comme les données de marché, les passerelles d'ordres, la plateforme de recherche ou le déploiement, et est mesurée sur leur fiabilité et sur la
rapidité de livraison des équipes produit.
\end{definition}

Le \omterm{def:fm:technology-strategy:tier}{palier de latence} traverse les deux. L'infrastructure du palier~3 est en général une plateforme, puisque les plateformes de négociation, les baies et les routes sont partagées par toutes les stratégies
qui courent~; le tutoriel n'attribue un palier propre à chaque stratégie que par souci de clarté. Ce qui relève d'une plateforme, c'est ce dont plusieurs \omterm{def:fm:technology-strategy:teams}{équipes produit} ont besoin
sous la même forme~; ce qui relève d'une \omterm{def:fm:technology-strategy:teams}{équipe produit}, c'est ce qui rend une stratégie différente. Le \emph{Team Topologies} de Skelton et Pais (2019) fournit le
vocabulaire~; le chapitre~21 l'applique aux ingénieurs d'une société de trading (\cref{fig:fm:technology-strategy:teams}).

\begin{omfigure}
\begin{tikzpicture}[font=\footnotesize,>=stealth,box/.style={draw,rounded corners,align=center,minimum height=0.75cm,text width=2.3cm}]
\foreach \x/\n in {-4.2/MM futures,-1.4/arb ETF,1.4/MM options,4.2/arb stat} {\node[box,fill=omDef!15] at (\x,1.5) {\n\\{\scriptsize équipe produit}};}
\node[box,fill=omThm!15,text width=7.6cm] (lat) at (-1.4,0) {plateforme à faible latence~: plateformes de négociation, baies, routes, passerelles d'ordres (palier~3)};
\node[box,fill=gray!15,text width=10.6cm] (base) at (0,-1.1) {plateforme de base~: données de marché, déploiement, recherche, supervision};
\draw[->] (-4.2,0.4) -- (-4.2,1.1); \draw[->] (-1.4,0.4) -- (-1.4,1.1); \draw[->] (1.4,0.4) -- (1.4,1.1);
\draw[->] (4.2,-0.7) -- (4.2,1.1);
\end{tikzpicture}

\omcaption{\omterm{def:fm:technology-strategy:teams}{Équipes plateforme} et \omterm{def:fm:technology-strategy:teams}{équipes produit} de la firme du chapitre~: trois stratégies qui courent partagent une plateforme à faible latence~; l'équipe d'arbitrage statistique
n'utilise que la plateforme de base. Schéma.}\label{fig:fm:technology-strategy:teams}
\end{omfigure}

\section{Le budget technologique}

\begin{definition}[Dépenses de fonctionnement et de transformation]\label{def:fm:technology-strategy:runchange}
Les \emph{dépenses de fonctionnement}\index{dépenses de fonctionnement} (run the bank) sont les dépenses technologiques nécessaires pour maintenir les activités existantes en service~: infrastructure,
licences, frais de plateformes et de données, support de production et maintenance. Les \emph{dépenses de transformation}\index{dépenses de transformation} (change the bank) sont les dépenses
consacrées à de nouvelles capacités~: nouvelles stratégies, nouvelles plateformes, nouveaux services et remplacements qui changent ce que la firme peut faire.
\end{definition}

Les paliers retenus coûtent 14.97~millions de dollars par an, et le support de production 1.2~million de dollars~: 16.2~millions de dollars de fonctionnement. Les projets de l'année, de nouvelles
connexions aux plateformes de négociation, une mise à niveau du feed handler sur FPGA et une plateforme de recherche, coûtent 4.0~millions de dollars de transformation~: le fonctionnement représente 80.2~\% du
budget. Les dépôts réglementaires de la firme de l'exemple montrent pourquoi la répartition compte~: un coût qui n'a pas baissé alors que le chiffre d'affaires reculait de 29~\% entre 2020 et 2023 relève du fonctionnement, et chaque décision de palier s'y ajoute.
Le rapport annuel de Flow Traders chiffre ses dépenses technologiques à 70.6~millions d'euros en 2025, contre 66.6~millions d'euros en 2024, à côté d'un produit net de trading
de 485.8~millions d'euros~: une autre ligne fixe dans une activité variable.

\begin{method}[Définir la stratégie technologique]\label{met:fm:technology-strategy:set}
\begin{enumerate}
\item Pour chaque stratégie, estimer la part du revenu décidée par les courses et par les paliers des concurrents.
\item Chiffrer chaque palier à partir des frais publiés et des coûts propres à la firme~; choisir le palier qui maximise le profit de la stratégie, et consigner le choix
comme un engagement stratégique assorti d'une date de réexamen.
\item Mutualiser l'infrastructure de course sous forme de plateforme~; garder le code propre à chaque stratégie dans les \omterm{def:fm:technology-strategy:teams}{équipes produit}.
\item Présenter le budget en fonctionnement et en transformation, et exiger que chaque projet de transformation indique le coût de fonctionnement qu'il ajoutera.
\item Revoir le choix lorsque les concurrents bougent~: la course aux armements modifie la réponse chaque année.
\end{enumerate}
\end{method}

\section{Construction~: paliers de latence}\label{bld:fm:technology-strategy:techtier}

\begin{build}
\bfield{Rôle} Le côté stratégique de la vitesse~: coûts des paliers, gain de chaque palier stratégie par stratégie, équilibre de la course aux armements, et répartition
fonctionnement/transformation.
\bfield{Interface} \texttt{firm.techtier}~: \texttt{Tier}, \texttt{colobill\_venue\_fees} (sur \texttt{firm.colobill} du livre~14), \texttt{Strategy},
\texttt{revenue}, \texttt{best\_tier}, \texttt{arms\_race}, \texttt{waste}, \texttt{run\_change}.
\bfield{Règles} Les revenus font la somme du fonds~; le palier le plus rapide présent prend toutes les courses, partagées également~; un équilibre est un profil dont aucune firme ne change~;
lorsque les meilleures réponses cyclent, une erreur est levée au lieu de renvoyer un profil.
\bfield{Tests d'acceptation} \texttt{code/firm/techtier/tests/}~: l'identité du fonds~; l'absence de déviation profitable à l'équilibre~; les frais de plateforme issus d'un barème
publié~; la répartition du budget.
\bfield{Pour aller plus loin} La base de coûts de connectivité du livre~14 comme source des coûts de palier~; coûts de plateforme partagés entre stratégies~; équilibres mixtes lorsque les
meilleures réponses cyclent.
\end{build}

\begin{omsources}
\item E.~Budish, P.~Cramton et J.~Shim, ``The high-frequency trading arms race'', \emph{Quarterly Journal of Economics} 130(4), 2015~;
M.~Aquilina, E.~Budish et P.~O'Neill, ``Quantifying the high-frequency trading `arms race''', \emph{Quarterly Journal of Economics} 137(1), 2022.
\item Virtu Financial, Forms 10-K 2020--2025~; Flow Traders, Annual Report 2025.
\item M.~Skelton et M.~Pais, \emph{Team Topologies}, IT Revolution, 2019.
\item Barèmes de frais tels qu'ils figurent dans le livre~14, chapitre~9.
\end{omsources}

\section{Exercices}

\begin{exercise}[$\star$]\label{exo:fm:technology-strategy:1}
Calculer de quel pourcentage les dépenses de communication et de données de Virtu ont augmenté entre 2019 et 2025, et quelle part du chiffre d'affaires elles représentaient chaque année.
\end{exercise}

\begin{exercise}[$\star$]\label{exo:fm:technology-strategy:2}
Une stratégie a $V=60$, $\alpha=0.8$ et cinq sociétés. Calculer le gain d'une société dans le palier rapide seule, puis lorsque deux autres sociétés y sont aussi présentes.
\end{exercise}

\begin{exercise}[$\star$]\label{exo:fm:technology-strategy:3}
Calculer les frais de plateforme des trois paliers, et la part qu'ils représentent dans le coût de chaque palier.
\end{exercise}

\begin{exercise}[$\star\star$]\label{exo:fm:technology-strategy:4}
Utiliser la \cref{prop:fm:technology-strategy:count} pour prévoir combien de sociétés participent à la course dans chaque stratégie.
\end{exercise}

\begin{exercise}[$\star\star$]\label{exo:fm:technology-strategy:5}
Expliquer pourquoi la société ne devrait pas acheter le palier 3 pour l'arbitrage statistique alors qu'elle l'achète pour trois autres stratégies.
\end{exercise}

\begin{exercise}[$\star\star$]\label{exo:fm:technology-strategy:6}
Classer en run ou en change~: renouveler le contrat d'une baie, construire une nouvelle passerelle vers une plateforme, remplacer un commutateur en fin de vie, la facture des données de marché.
\end{exercise}

\begin{exercise}[$\star\star\star$]\label{exo:fm:technology-strategy:7}
\emph{Programmation.} Doubler les autres coûts du palier 3 pour les porter à 7~millions de dollars et relancer \texttt{fm\_tech.races()}. Déterminer ce qu'il advient des équilibres et du gaspillage.
\end{exercise}

\begin{exercise}[$\star\star\star$]\label{exo:fm:technology-strategy:8}
\emph{Trouver la faille.} «~Les dépenses liées à la vitesse sont gaspillées, comme le montrent les modèles~; nous arrêterons les mises à niveau et économiserons l'argent.~»
\end{exercise}

\section{Problème~: le palier le plus rapide}

\begin{problem}[{Problème du week-end --- le palier le plus rapide}]\label{pb:fm:technology-strategy:1}
Un directeur technique doit proposer dans quelles stratégies la société participera à la course, et à quelle vitesse, pour les trois prochaines années.

\medskip\noindent\textbf{Partie I --- La vitesse comme stratégie.}
\begin{enumerate}
\item Définir un \omterm{def:fm:technology-strategy:tier}{palier de latence} et une \omterm{def:fm:technology-strategy:race}{course aux armements de la latence}.
\item Expliquer ce que révèlent les dépenses de communication et de données de Virtu.
\item Dire ce qu'Aquilina, Budish et O'Neill ont établi sur les courses.
\item Chiffrer les trois paliers et préciser d'où viennent les prix.
\end{enumerate}

\medskip\noindent\textbf{Partie II --- Le choix.}
\begin{enumerate}[resume]
\item Écrire le modèle de course.
\item Donner le profit par palier pour chaque stratégie, ainsi que le meilleur palier.
\item Expliquer pourquoi la réponse diffère pour le basis trade crypto.
\item Dire ce qui modifierait la réponse pour l'arbitrage statistique.
\end{enumerate}

\medskip\noindent\textbf{Partie III --- La course.}
\begin{enumerate}[resume]
\item Énoncer et démontrer la \cref{prop:fm:technology-strategy:count}.
\item Donner le profil d'équilibre et les dépenses de chaque stratégie.
\item Donner la part gaspillée par stratégie et au total.
\item Expliquer pourquoi le gaspillage du secteur n'est pas celui de la société.
\item Dire ce que concluent Budish, Cramton et Shim.
\end{enumerate}

\medskip\noindent\textbf{Partie IV --- L'organisation et le budget.}
\begin{enumerate}[resume]
\item Définir les \omterm{def:fm:technology-strategy:teams}{équipes plateforme} et les \omterm{def:fm:technology-strategy:teams}{équipes produit}, et les schématiser pour la société.
\item Définir les \omterm{def:fm:technology-strategy:runchange}{dépenses de fonctionnement} (run the bank) et les \omterm{def:fm:technology-strategy:runchange}{dépenses de transformation} (change the bank).
\item Donner la répartition entre run et change de la société.
\item Expliquer ce que révèle le rapport de Flow Traders.
\item Dire à quelle fréquence le choix du palier doit être réexaminé, et sur quels éléments.
\item Énoncer le \emph{résultat nommé}~: le palier qui maximise le profit pour chacune des cinq stratégies, et la part des dépenses liées à la vitesse du secteur que
l'équilibre de la course aux armements gaspille par rapport à une situation où personne ne se mettrait à niveau.
\item En deux phrases, rédiger la stratégie technologique.
\end{enumerate}
\end{problem}

\section{Questions d'entretien}

\begin{interviewq}[$\star$ \iqroles{développeur}]\label{iq:fm:technology-strategy:1}
Quelle est la différence entre une \omterm{def:fm:technology-strategy:teams}{équipe plateforme} et une \omterm{def:fm:technology-strategy:teams}{équipe produit} dans une société de trading~?
\end{interviewq}

\begin{interviewq}[$\star$ \iqroles{trader}]\label{iq:fm:technology-strategy:2}
Pourquoi une société peut-elle perdre de l'argent en étant plus rapide~?
\end{interviewq}

\begin{interviewq}[$\star\star$ \iqroles{chercheur}]\label{iq:fm:technology-strategy:3}
Cinq sociétés se disputent un pot de 30~millions de dollars, dont 90~\% sont attribués par des courses~; être rapide coûte 4~millions de dollars de plus par an. Combien d'entre elles seront rapides~?
\end{interviewq}

\begin{interviewq}[$\star\star$ \iqroles{développeur, trader}]\label{iq:fm:technology-strategy:4}
Comment décideriez-vous si une stratégie a besoin d'un FPGA~?
\end{interviewq}

\begin{interviewq}[$\star\star$ \iqroles{risque}]\label{iq:fm:technology-strategy:5}
Pourquoi la majeure partie du budget technologique d'une société de trading devient-elle des \omterm{def:fm:technology-strategy:runchange}{dépenses de fonctionnement} (run the bank)~?
\end{interviewq}

\begin{interviewq}[$\star\star\star$ \iqroles{chercheur}]\label{iq:fm:technology-strategy:6}
Les concepteurs de marchés proposent des enchères par lots fréquentes. Que feraient-ils de la course aux armements dans le modèle du chapitre~?
\end{interviewq}
